\chapter{Conclusion}

The goal of this thesis was to implement Field Oriented Control firmware on STM-32 type microcontroller for an existing uFOC motor driver 
developed by the Department of Control Engineering at the Czech Technical University in Prague.

The high level goal of the uFOC project is to develop a compact integrated servo actuator inspired by the modular actuator architecture pioneered by Benjamin Katz at MIT. \cite{ben_katz}
The concept combines motor driver electronics, sensing, communication interfaces and closed loop motor control into a compact device directly mountable onto the motor assembly. 
Such integration reduces external wiring complexity and enables development of low cost modular actuators suitable for robotics and other mechatronic applications.

The implemented firmware provides software modules for communication with peripherals such as the DRV8316C gate driver and the MT6835 magnetic encoder, 
as well as torque, velocity and position control loops based on the Field Oriented Control principle with emphasis on deterministic behaviour.

The firmware also implements CAN bus communication together with a client side Python interface.
The communication protocol supports device identification, allowing multiple drivers to share a single CAN bus line in a daisy chain configuration. 
The implemented command interface allows modification of control targets, PI and PID controller constants, 
software limits and reading of processed sensor data and runtime state variables.

The first part of this thesis provided a brief overview of the theoretical principles required for implementation of FOC control, 
including Clarke and Park transformations, Space Vector PWM modulation and cascaded current, velocity and position control loops.

The second part of this thesis described the practical firmware implementation in detail, including firmware architecture, MCU and peripheral configuration, 
encoder processing, current sensing, implementation of PI and PID regulators, practical realization of the FOC algorithm and CAN communication handling.

Experimental testing confirmed correct operation of the implemented torque, velocity and position control loops on the iPower GM4108 motor. 
The firmware demonstrated stable torque regulation, accurate velocity tracking and smooth position control. 
Communication between multiple drivers connected through a shared CAN bus was also successfully demonstrated.

The developed firmware fulfills the goals defined in the thesis assignment and provides a functional embedded
FOC control platform suitable for further development and experimental applications. 
Several future improvements were proposed, including optimization of the torque control loop execution time, 
encoder calibration improvements, anticogging compensation, synchronized multi-axis control and IMU integration.
